\documentclass[a4paper, 11pt]{article}
\usepackage[utf8]{inputenc}
\usepackage[T1]{fontenc}
\usepackage{lmodern}
\usepackage[margin=1in]{geometry}
\usepackage{amsmath, amssymb}
\usepackage{booktabs}
\usepackage[protrusion=true,expansion=false]{microtype}
\usepackage{xcolor}
\usepackage{listings}
\usepackage{tcolorbox}
\usepackage{url}
\usepackage{graphicx}

\definecolor{codegreen}{rgb}{0.15, 0.55, 0.15}
\definecolor{codegray}{rgb}{0.45, 0.45, 0.45}
\definecolor{codeblue}{rgb}{0.12, 0.35, 0.75}
\definecolor{backcolour}{rgb}{0.96, 0.97, 0.98}
\definecolor{termbg}{rgb}{0.14, 0.15, 0.17}
\definecolor{termfg}{rgb}{0.92, 0.92, 0.92}

\lstdefinestyle{erlstyle}{
    backgroundcolor=\color{backcolour},   
    commentstyle=\color{codegreen}\itshape,
    keywordstyle=\color{codeblue}\bfseries,
    numberstyle=\tiny\color{codegray},
    stringstyle=\color{codegray},
    basicstyle=\ttfamily\footnotesize,
    breaklines=true,                 
    keepspaces=true,                 
    tabsize=2,
    frame=single,
    rulecolor=\color{black!10}
}
\lstset{style=erlstyle}

\tcbset{
    terminalbox/.style={
        colback=termbg,
        colupper=termfg,
        colframe=black!65,
        arc=3pt,
        boxrule=1pt,
        fontupper=\ttfamily\footnotesize,
        left=6pt, right=6pt, top=5pt, bottom=5pt
    }
}

\title{\textbf{Chordy: A Distributed Hash Table}}
\author{\textbf{Rodrigo Perestrelo, Lennard Last} \\ ID2201 Distributed Systems, Basic Course}
\date{\today{}}

\begin{document}

\maketitle

\section{Introduction}

The objective of this assignment is to design, implement, and empirically evaluate \textbf{Chordy}, a structured Distributed Hash Table based on the Chord protocol scheme. The system organizes autonomous Erlang processes into an overlay identifier ring where keys and nodes share a uniform discrete space.

\section{Iterative Architectural Progression}

We developed the system through two foundational stages, transitioning from topological ring maintenance to a fully functioning distributed storage engine.

\subsection{Iteration 1: Topology and Self-Stabilizing Ring (\texttt{node1})}
The baseline architecture, implemented in modules \texttt{key} and \texttt{node1}, maintains a circular overlay network without data storage.
\begin{itemize}
    \item \textbf{Circular Identifier Space:} Nodes generate 30-bit random identifiers ($1$ to $1\,000\,000\,000$). The predicate \texttt{key:between(Key, From, To)} resolves circular intervals $(From, To]$ by handling standard boundaries ($From < To$), wraparounds ($From > To$), and the full-circle degenerate identity ($From == To$).
    \item \textbf{Asynchronous State Maintenance:} Each node tracks its current identifier, a predecessor pointer, and a successor pointer as $\{Key, Pid\}$ tuples. Ring integrity is maintained asynchronously via four core message primitives:
    \begin{itemize}
        \item \texttt{\{key, Qref, Peer\}}: Resolves the key of a bootstrap node during connection.
        \item \texttt{\{request, Peer\}}: Solicits predecessor information from the successor.
        \item \texttt{\{status, Pred\}}: Transmits predecessor status to the querying predecessor.
        \item \texttt{\{notify, New\}}: Proposes a friendly candidate as a potential new predecessor.
    \end{itemize}
    \item \textbf{Periodic Stabilization:} Driven by \texttt{timer:send\_interval/3}, each node queries its successor's predecessor via \texttt{request}. If an intermediate node $X$ is discovered such that $X \in (Id, Successor]$, the querying node adopts $X$ as its successor. It then informs the successor of its presence via \texttt{notify}.
\end{itemize}

\begin{lstlisting}[language=erlang]
stabilize(Pred, Id, Successor) ->
    {Skey, Spid} = Successor,
    case Pred of
        nil ->
            Spid ! {notify, {Id, self()}},
            Successor;
        {Id, _} -> Successor;
        {Skey, _} ->
            Spid ! {notify, {Id, self()}},
            Successor;
        {Xkey, Xpid} ->
            case key:between(Xkey, Id, Skey) of
                true ->
                    Xpid ! {request, self()},
                    {Xkey, Xpid};
                false ->
                    Spid ! {notify, {Id, self()}},
                    Successor
            end
    end.
\end{lstlisting}

\noindent \textbf{Structural Insight:} When a node points to itself ($Predecessor == Id$), the condition $From == To$ triggers true in \texttt{key:between/3}, ensuring that incoming nodes are accepted into single-node rings without explicit corner-case branches.

\subsection{Iteration 2: Distributed Store and Dynamic Handover (\texttt{node2})}
Module \texttt{node2} integrates a local key-value repository, delegating key-range ownership across ring participants:
\begin{itemize}
    \item \textbf{Responsibility Invariant:} A node $Id$ with predecessor $Pkey$ owns the half-open interval $(Pkey, Id]$. Operations on keys falling outside this interval are forwarded sequentially to the immediate successor \texttt{Spid}.
    \item \textbf{Transparent Client Routing:} Client queries (\texttt{\{add, Key, Value, Qref, Client\}} and \texttt{\{lookup, Key, Qref, Client\}}) carry a unique tag \texttt{Qref}. When forwarded along the ring, the owning node sends the completion response directly to \texttt{Client}, eliminating intermediate return hops.
    \item \textbf{Dynamic Storage Delegation (\texttt{handover}):} When a new predecessor $Nkey$ joins, the interval $(Pkey, Id]$ is partitioned. The local node retains $(Nkey, Id]$ and delegates $(Pkey, Nkey]$ by dispatching a \texttt{\{handover, Rest\}} message to $Npid$.
\end{itemize}

\begin{lstlisting}[language=erlang]
handover(Id, Store, Nkey, Npid) ->
    {Keep, Rest} = storage:split(Nkey, Id, Store),
    Npid ! {handover, Rest},
    Keep.
\end{lstlisting}

\section{Empirical Evaluation}

\subsection{Evaluation 1: Single-Node and Multi-Node Probe Verification}
We initialized a distributed ring across separate Erlang VM nodes to measure ring transit latency and verify pointer convergence using a circumnavigating probe message:
\begin{tcolorbox}[terminalbox]
(first@127.0.0.1)1> P1 = node1:start(100). \\
<0.94.0> \\
(first@127.0.0.1)2> register(chord\_root, P1). \\
true \\
(first@127.0.0.1)3> chord\_root ! probe. \\
probe \\
Probe returned in 419 microseconds after 3 hops: [<10738.94.0>, <10737.94.0>, <0.94.0>]
\end{tcolorbox}
\textbf{Observed Result:} The probe completed a full circle across 3 independent OS-level Erlang processes in 419 microseconds, confirming that periodic stabilization converged into a consistent, closed ring topology.

\subsection{Evaluation 2: Scalability and Performance Profiling}
We executed sequential and concurrent benchmark routines to evaluate the system under varying ring topologies ($N=1$, $N=2$, and $N=4$ nodes) storing $1\,000$ elements:

\begin{tcolorbox}[terminalbox]
8> benchmark:benchmark(1000, N1).
Starting addition of 1000 elements...
Addition completed in 6361 us (6 ms)
Starting lookup of 1000 elements...
Lookup completed in 1818 us (1 ms)
{6361, 1818}

11> benchmark:benchmark(1000, N1).
Starting addition of 1000 elements...
Addition completed in 9713 us (9 ms)
Starting lookup of 1000 elements...
Lookup completed in 2430 us (2 ms)
{9713, 2430}

15> benchmark:benchmark(1000, N1).
Starting addition of 1000 elements...
Addition completed in 8413 us (8 ms)
Starting lookup of 1000 elements...
Lookup completed in 2697 us (2 ms)
{8413, 2697}

16> benchmark:benchmark(1000, N1, 4).
Starting benchmark with 4 clients (250 keys each)...
Parallel benchmark completed in 12469 us (12 ms)
12469
\end{tcolorbox}

\begin{table}[h]
\centering
\small
\begin{tabular}{@{}lcccc@{}}
\toprule
\textbf{Configuration} & \textbf{Add Time ($\mu$s)} & \textbf{Lookup Time ($\mu$s)} & \textbf{Avg Add/Op} & \textbf{Avg Lookup/Op} \\ \midrule
1 Node ($1\,000$ keys)  & $6\,361$  & $1\,818$ & $6.36\,\mu$s & $1.82\,\mu$s \\
2 Nodes ($1\,000$ keys) & $9\,713$  & $2\,430$ & $9.71\,\mu$s & $2.43\,\mu$s \\
4 Nodes ($1\,000$ keys) & $8\,413$  & $2\,697$ & $8.41\,\mu$s & $2.70\,\mu$s \\
4 Concurrent Clients    & \multicolumn{2}{c}{$12\,469\,\mu$s (Total Pipeline)} & \multicolumn{2}{c}{Contention on $N_1$ Mailbox} \\ \bottomrule
\end{tabular}
\caption{Empirical benchmark timings across ring configurations.}
\end{table}

\subsection{Evaluation 3: Architectural Analysis and Limiting Factors}
\begin{itemize}
    \item \textbf{Single-Node vs Multi-Node Overhead:} In a single-node configuration, $100\%$ of requests are local, eliminating inter-process message passing. When scaling to 2 nodes, approximately $50\%$ of operations require forwarding to the successor, increasing addition latency from $6\,361\,\mu$s to $9\,713\,\mu$s.
    \item \textbf{Local Complexity Mitigation ($N=4$):} Expanding the ring to 4 nodes reduced addition latency to $8\,413\,\mu$s. Because each node stores approximately $250$ keys instead of $1\,000$, the $O(K)$ traversal cost of \texttt{lists:keystore/4} drops significantly, offsetting message routing overhead.
    \item \textbf{Successor-Hop Forwarding Penalty:} Lookup latency rose monotonically ($1\,818\,\mu\text{s} \rightarrow 2\,430\,\mu\text{s} \rightarrow 2\,697\,\mu\text{s}$). Because our ring routes via immediate successors ($O(N)$ hops) without finger routing tables ($O(\log N)$), each additional node lengthens the traversal path.
    \item \textbf{The Mailbox Bottleneck:} When 4 client processes submit requests concurrently through a single entrypoint ($N_1$), latency climbs to $12\,469\,\mu$s. The entry node's mailbox serializes all inbound traffic, creating an architectural bottleneck.
\end{itemize}

\section{Handling of Node Failure}
In the third iteration of \textbf{Chordy}, the handling of failures and subsequent self-healing was added. The implementation has the following three key additions.
\subsection{Neighbor Monitoring}
Each node now monitors its neighbors using Erlang's built-in \texttt{monitor} mechanism. This way a node will be notified when its predecessor or successor fails by receiving a \texttt{\{'DOWN', Ref, ..\}} message and can react accordingly.
\subsection{Next Pointer}
During stabilization a node will now also pass its own successor in the request reply and the recipient of the status message will save it in the \texttt{Next} variable of the node. This pointer is needed because a node that loses its successor would otherwise have no way to reach the rest of the ring. When a closer successor is found during stabilization, the old successor becomes the new \texttt{Next}.
\begin{lstlisting}[language=erlang]
request(Peer, Predecessor, {Skey, _, Spid}) ->
    case Predecessor of
        nil ->
            Peer ! {status, nil, {Skey, Spid}};
        {Pkey, _, Ppid} ->
            Peer ! {status, {Pkey, Ppid}, {Skey, Spid}}
    end.
\end{lstlisting}
\begin{lstlisting}[language=erlang]
node(Id, Predecessor, Successor, Next, Store) ->
    receive
        ...
        {status, Pred, Nx} ->
            {Succ, Nxt} = stabilize(Pred, Nx, Id, Successor),
            node(Id, Predecessor, Succ, Nxt, Store);
\end{lstlisting}
\subsection{Failure Handling}
When the successor fails, the node will set \texttt{Next} as its new successor and trigger stabilization to inform its new successor and receive a new \texttt{Next} pointer. When the predecessor fails, the node will just set it to \texttt{nil} and wait for the new predecessor's stabilization.
\begin{lstlisting}[language=erlang]
down(Ref, {_, Ref, _}, Successor, Next) ->
    %% Predecessor died
    {nil, Successor, Next};
down(Ref, Predecessor, {_, Ref, _}, {Nkey, Npid}) ->
    %% Successor died
    drop(Ref),
    Nref = monitor(Npid),
    self() ! stabilize,
    {Predecessor, {Nkey, Nref, Npid}, nil}.
\end{lstlisting}

\section{Replication}
The self-healing capabilities in the third iteration may keep the ring up and running, but a node that fails takes all its data with it. To counter this a simple form of replication is implemented in the fourth iteration.

Now every node will store a \texttt{Replica} of its predecessor's store. This is a sensible design decision, since upon exit of a node, its successor is now responsible for its data and thus can just move \texttt{Replica} into its own \texttt{Store}.
\begin{lstlisting}[language=erlang]
down(Ref, {_, Ref, _}, Successor, Next, Store, Replica) ->
    %% Predecessor died
    MergedStore = storage:merge(Store, Replica),
    NewReplica = storage:create(),
    {nil, Successor, Next, MergedStore, NewReplica};
\end{lstlisting}
When a node adds a key to its own store, it forwards the pair to its successor in a \texttt{replicate} message. The successor stores it in its \texttt{Replica} and sends the confirmation to the client, so the client only receives \texttt{ok} once both copies exist.
\begin{lstlisting}[language=erlang]
{replicate, Key, Value, Qref, Client} ->
    NewReplica = storage:add(Key, Value, Replica),
    Client ! {Qref, ok},
    node(Id, Predecessor, Successor, Next, Store, NewReplica);
\end{lstlisting}
To make sure replication stays consistent through changes of the ring, the \texttt{handover} function now also returns the part of the storage that is the responsibility of the new predecessor and the \texttt{notify} function then returns it to be set as the new \texttt{Replica}.
\begin{lstlisting}[language=erlang]
handover(Id, Store, Nkey, Npid) ->
    {Keep, Rest} = storage:split(Nkey, Id, Store),
    Npid ! {handover, Rest},
    {Keep, Rest}.
\end{lstlisting}
\begin{lstlisting}[language=erlang]
{notify, New} ->
    case notify(New, Id, Predecessor, Store) of
        {Pred, NewStore, nil} ->
            node(Id, Pred, Successor, Next, NewStore, Replica);
        {Pred, NewStore, NewReplica} ->
            node(Id, Pred, Successor, Next, NewStore, NewReplica)
    end;
\end{lstlisting}

\section{Conclusions}

This laboratory demonstrates the mechanics of structured decentralized storage over an identifier ring. Decentralized ring stabilization allows autonomous nodes to dynamically self-organize without a centralized coordinator. While distributed partitioning divides local storage costs, linear successor routing incurs an $O(N)$ hop penalty, and single-entry client access introduces mailbox serialization. Achieving high-throughput scalability requires finger-table routing and multi-point client dispatching.

\end{document}